% 5.1 章 · 无约束优化与拟牛顿
% 这一章的笔记要把"让 f 变小"拆成四件互不相同的事：
%   1) 步长策略（定步长 / 线搜索 / 信赖域）决定"能走多远"；
%   2) 曲率信息（有没有、准不准）决定"走多快"；
%   3) 收敛的**形态**（线性 / 超线性 / 二次）与收敛的**速度**（迭代数）是两码事；
%   4) 停机判据是结论的一部分：换一个判据就换一个结论。
% 所有结论都配一条带实测数字的断言，附表由 tools/build_book.py 生成后 \input 进来。

\numlabchapter{无约束优化与拟牛顿}{5.1}

\section{这一章要解决什么}

给定 $f:\mathbb{R}^n\to\mathbb{R}$，求 $\min_x f(x)$。教科书在这里通常会给出一串方法：
梯度下降、牛顿法、拟牛顿、共轭梯度、信赖域，然后说"牛顿法二次收敛、拟牛顿超线性收敛"。
本章要做的不是把这张表抄一遍，而是把**这些话成立的条件与代价**算清楚：

\begin{enumerate}
  \item \textbf{梯度下降慢在哪里，慢多少？} 慢在条件数 $\kappa=L/\mu$：
        迭代数是 $\Theta(\kappa\ln(1/\varepsilon))$，这条结论可以逐位验证（实测
        $\kappa=1000,\varepsilon=10^{-8}$ 时 $9211$ 步，理论值 $9211$ 步）。
  \item \textbf{"二次收敛"到底是什么条件下的事？} 它要求"解附近 $+$ Hessian 正定 $+$
        \textbf{整步被接受} $+$ 已经在二次区"。缺一条就退化：卡住步长后退化为线性，
        线性因子实测就是 $1-\alpha_{\text{cap}}$。
  \item \textbf{没有二阶导数时，割线信息能替代多少？} BFGS 的割线方程是**更新公式的
        目标**而不是恒等式：更新后残差 $3\times10^{-15}$，更新前 $0.999$，差 12 个数量级。
        L-BFGS 只精确满足最近 $m$ 对，$m=2$ 时最新一对 $0.0$、上一对 $0.13$。
  \item \textbf{鞍点上会发生什么？} 梯度下降在稳定流形上\textbf{收敛到鞍点并停住}；
        扰动只有落在负曲率方向 $u_{\min}$ 上才有用，逃逸阈值是
        $\delta^\star=R(1+|\lambda_{\min}|\alpha)^{-N}$，实测
        $2.866\times10^{-3}$ 对解析值 $(1.05)^{-120}=2.866\times10^{-3}$。
\end{enumerate}

记号：$u=\tfrac12\varepsilon_{\text{mach}}=2^{-53}\approx1.11\times10^{-16}$ 是单位舍入；
$\kappa=L/\mu$ 是 Hessian 的条件数（二次型上 $L=\lambda_{\max}$、$\mu=\lambda_{\min}$）；
$e_k=x_k-x^\star$；$\|\cdot\|$ 缺省为二范数。

\section{二次型上的梯度下降}

本章所有"速率"结论都先在严格凸二次型上说清楚，因为只有在那里每一步都能算到底。

\begin{theorem}[定步长梯度下降的误差递推与收敛充要条件]
\label{thm:gd-recurrence}
设 $f(x)=\tfrac12x^{\mathsf T}Ax-b^{\mathsf T}x$，$A=A^{\mathsf T}$，
$0<\mu I\preceq A\preceq LI$，$x^\star=A^{-1}b$。
迭代 $x_{k+1}=x_k-\alpha\nabla f(x_k)$ 的误差满足
\[
  e_{k+1}=(I-\alpha A)e_k ,
\]
并且对任意初值收敛的充要条件是
\[
  0<\alpha<\frac{2}{L}.
\]
此时 $\|e_k\|_2\le\rho(\alpha)^k\|e_0\|_2$，
$\rho(\alpha)=\max_i|1-\alpha\lambda_i|=\max\{|1-\alpha\mu|,\,|1-\alpha L|\}$。
\end{theorem}

\begin{proof}
$\nabla f(x)=Ax-b=A(x-x^\star)=Ae$，代入迭代式得
$e_{k+1}=e_k-\alpha Ae_k=(I-\alpha A)e_k$。
$A$ 对称，取正交特征基 $A=Q\Lambda Q^{\mathsf T}$，令 $c_k=Q^{\mathsf T}e_k$，
则 $c_{k+1}=(I-\alpha\Lambda)c_k$，逐分量 $c_{k+1}^{(i)}=(1-\alpha\lambda_i)c_k^{(i)}$，
于是 $c_k^{(i)}=(1-\alpha\lambda_i)^k c_0^{(i)}$，取二范数即得
$\|e_k\|_2\le\max_i|1-\alpha\lambda_i|^k\|e_0\|_2$。

\emph{充要性}：若 $0<\alpha<2/L$，则对每个 $i$ 有
$-1<1-\alpha L\le1-\alpha\lambda_i\le1-\alpha\mu<1$，故 $\rho<1$，迭代收敛。
反之若 $\alpha\ge2/L$，则 $|1-\alpha L|\ge1$；取 $e_0$ 为 $L$ 对应的特征向量，
$\|e_k\|=\rho^k\|e_0\|$ 不趋于零；$\alpha\le0$ 时 $|1-\alpha\lambda_i|\ge1$ 同样不收敛。
\end{proof}

\begin{corollary}[最优步长与依赖条件数的迭代数]
\label{cor:gd-rate}
$\rho(\alpha)$ 在
\[
  \alpha^\star=\frac{2}{L+\mu}
\]
处取到最小值 $\rho^\star=\dfrac{L-\mu}{L+\mu}=\dfrac{\kappa-1}{\kappa+1}$。
若要求 $\|e_N\|/\|e_0\|\le\varepsilon$，则
\[
  N(\kappa,\varepsilon)=\left\lceil\frac{\ln\varepsilon}{\ln\rho^\star}\right\rceil
  =\Theta\!\left(\kappa\ln\frac1\varepsilon\right),
  \qquad
  N\approx\frac{\kappa+1}{2}\ln\frac1\varepsilon .
\]
\end{corollary}

\begin{proof}
对 $\alpha\le\alpha^\star$ 有 $\alpha(L+\mu)\le2$，即 $1-\alpha\mu\ge\alpha L-1$，
故 $\rho(\alpha)=1-\alpha\mu$ 随 $\alpha$ 递减；对 $\alpha\ge\alpha^\star$ 同理得
$\rho(\alpha)=\alpha L-1$ 递增。所以最小值在 $\alpha^\star$ 处，值为
$1-\alpha^\star\mu=1-\frac{2\mu}{L+\mu}=\frac{L-\mu}{L+\mu}$。

迭代数：$\|e_N\|\le(\rho^\star)^N\|e_0\|$ 取等号当 $e_0$ 是 $\mu$ 的特征向量，
故所需步数是 $\lceil\ln\varepsilon/\ln\rho^\star\rceil$（该特征分量决定最后一段）。
最后用 $\ln\rho^\star=\ln\frac{\kappa-1}{\kappa+1}
=\ln\bigl(1-\frac{2}{\kappa+1}\bigr)\approx-\frac{2}{\kappa+1}$
得到 $N\approx\frac{\kappa+1}{2}\ln\frac1\varepsilon$。
实测（\texttt{test\_gd\_iteration\_count\_is\_linear\_in\_kappa}）：$\kappa=10,30,100,300,1000,3000,10000$
上的实测迭代数是 $92,277,922,2764,9211,27632,92104$，与
$\lceil\ln\varepsilon/\ln\rho^\star\rceil$ 的相对差小于 $10^{-3}$，
幂律拟合的指数为 $1.000$，常数 $N/\kappa$ 与 $\tfrac12\ln(1/\varepsilon)$ 之比在 $0.02$ 之内。
\end{proof}

\begin{remark}[锯齿的几何]
$e_{k+1}$ 的每一个特征分量独立地乘上 $1-\alpha\lambda_i$。$\lambda_i=L$ 的分量每步翻号
（因为 $1-\alpha^\star L=-\frac{\kappa-1}{\kappa+1}$），$\lambda_i=\mu$ 的分量却几乎不变号，
两者交替占优就形成等高线上的锯齿。$\kappa$ 越大，$\rho^\star$ 越接近 $1$
（$\kappa=100$ 时 $0.980198$，$\kappa=10^4$ 时 $0.999800$，实测与理论逐位相同），
所以"病态问题上的最速下降等于没走"这句话的量化形式就是 $\rho^\star=1-2/\kappa+O(\kappa^{-2})$。
\end{remark}

\subsection{临界步长：$2/L$ 与 $2/(L+\mu)$ 不是同一件事}

$\alpha^\star=2/(L+\mu)$ 是最优步长，而收敛域的边界是 $2/L$。两者的比是
$\frac{2/L}{2/(L+\mu)}=1+\frac1\kappa$：$\kappa=100$ 时只差 $1\%$。
这解释了一个常见困惑——"我明明取了最优步长，为什么稍微调大一点就发散了"：
从 $\alpha^\star$ 出发只有 $1/\kappa$ 的相对余量。
实测（\texttt{test\_gd\_step\_threshold\_is\_bisected\_not\_copied}）：以"$2000$ 步内误差降 $10$ 倍"
为判据二分，得到临界步长 $\alpha_{\text{thr}}/\alpha^\star$ 与
$(1+0.1^{1/N})(L+\mu)/(2L)$ 的相对差小于 $5\times10^{-3}$，
即 $\alpha_{\text{thr}}\to2/L$。

\section{重球动量}

\begin{theorem}[重球动量的谱半径]
\label{thm:momentum}
迭代 $x_{k+1}=x_k-\alpha\nabla f(x_k)+\beta(x_k-x_{k-1})$ 在二次型上的误差满足
$e_{k+1}=(I-\alpha A)e_k+\beta(e_k-e_{k-1})$。取
\[
  \alpha=\frac{4}{(\sqrt L+\sqrt\mu)^2},\qquad
  \beta=\left(\frac{\sqrt\kappa-1}{\sqrt\kappa+1}\right)^2
  =\left(\frac{\sqrt L-\sqrt\mu}{\sqrt L+\sqrt\mu}\right)^2 ,
\]
则每个特征方向上的特征多项式
$z^2-(1+\beta-\alpha\lambda_i)z+\beta=0$ 的两根模长都等于 $\sqrt\beta$。
\end{theorem}

\begin{proof}
令 $c_i=1+\beta-\alpha\lambda_i$。记 $s=\sqrt L$、$m=\sqrt\mu$，则
$\beta=\bigl(\frac{s-m}{s+m}\bigr)^2$、$\alpha=4/(s+m)^2$、
$2\sqrt\beta=\frac{2(s-m)}{s+m}$。直接计算两个端点：
\[
  1+\beta-\alpha L=\frac{2(s^2+m^2)-4s^2}{(s+m)^2}
  =\frac{-2(s-m)(s+m)}{(s+m)^2}=-2\sqrt\beta,
\]
\[
  1+\beta-\alpha\mu=\frac{2(s^2+m^2)-4m^2}{(s+m)^2}
  =\frac{2(s-m)(s+m)}{(s+m)^2}=+2\sqrt\beta .
\]
$c_i$ 是 $\lambda_i$ 的仿射函数，故 $\lambda_i\in[\mu,L]$ 时
$|c_i|\le2\sqrt\beta$。
二次方程两根之积为 $\beta>0$：若判别式 $c_i^2-4\beta<0$，两根是共轭复根，
模长都是 $\sqrt\beta$；若 $c_i^2=4\beta$（只可能在端点取到），两根都是实数
$\pm\sqrt\beta$，模长仍为 $\sqrt\beta$。
\end{proof}

\begin{corollary}[速率是渐近的，迭代数不是]
\label{cor:momentum-transient}
由定理 \ref{thm:momentum}，每个特征分量的二阶递推矩阵的特征值模长都是
$\sqrt\beta=\frac{\sqrt\kappa-1}{\sqrt\kappa+1}$，故
$\|e_k\|\le C(1+k)\beta^{k/2}\|e_0\|$；即\textbf{渐近速率}是
$\frac{\sqrt\kappa-1}{\sqrt\kappa+1}$，但带一个 $O(k)$ 的瞬态因子。
因此"动量把迭代数降到 $1/\sqrt\kappa$"只是渐近速率的推论。
实测（\texttt{test\_momentum\_rate\_matches\_theory\_but\_iteration\_ratio\_does\_not}）：
$\kappa=100,1000,10^4$ 上实测速率与理论值的相对差只有 $0.009,0.002,0.0007$（速率对得上），
而迭代数之比 $N_{\text{gd}}/N_{\text{mom}}$ 是 $8.06,24.39,74.13$，
只有 $\sqrt\kappa(=10,31.6,100)$ 的 $0.81,0.77,0.74$ 倍。
把 $\sqrt\kappa$ 当成迭代数结论会高估 $20\%$ 到 $26\%$。
\end{corollary}

\begin{proof}
重根只出现在 $|c_i|=2\sqrt\beta$ 的两个端点上（$\lambda_i=\mu$ 或 $L$）。
此时二阶递推的转移矩阵是 Jordan 块，贡献一个 $k$ 的线性因子；
其余方向的转移矩阵可对角化，因子有界。合并得
$\|e_k\|\le C(1+k)(\sqrt\beta)^k\|e_0\|$。
$C$ 与 $k$ 无关，故渐近对数速率仍是 $\ln\sqrt\beta$；
但要在有限步数内达到给定精度，就必须多付这个 $(1+k)$ 因子，
这正是实测迭代数比 $1/\sqrt\kappa$ 差两成的原因。
\end{proof}

\section{牛顿法：它到底承诺了什么}

\begin{theorem}[二次型上一步到驻点]
\label{thm:newton-one-step}
若 $f$ 是定理 \ref{thm:gd-recurrence} 中的二次型且 $A$ 非奇异，
则从任意 $x_0$ 出发，整步牛顿法 $x_1=x_0-\nabla^2f(x_0)^{-1}\nabla f(x_0)$
满足 $x_1=x^\star$。
\end{theorem}

\begin{proof}
$\nabla^2f(x_0)=A$ 与 $x_0$ 无关，$\nabla f(x_0)=Ax_0-b$，
故 $x_1=x_0-A^{-1}(Ax_0-b)=A^{-1}b=x^\star$。
\end{proof}

\begin{theorem}[局部二次收敛及其常数]
\label{thm:newton-quadratic}
设 $f\in C^3$，$H=\nabla^2f(x^\star)\succ0$，$\mu=\lambda_{\min}(H)$，
且存在 $M>0$ 使 $\|\nabla^2f(x)-H\|\le M\|x-x^\star\|$（在 $x^\star$ 的邻域内）。
若第 $k$ 步的误差满足 $\|e_k\|\le r$ 且 $M\|e_k\|\le\mu/2$，则整步牛顿法给出
\[
  \|e_{k+1}\|\le\frac{3M}{\mu}\,\|e_k\|^2 .
\]
\end{theorem}

\begin{proof}
记 $H_k=\nabla^2f(x_k)$、$E_k=H_k-H$，由假设 $\|E_k\|\le M\|e_k\|$。
由中值形式的 Taylor 展开（$\nabla f(x^\star)=0$）：
\[
  \nabla f(x_k)=He_k+R_k,\qquad \|R_k\|\le\frac M2\|e_k\|^2 .
\]
代入 $x_{k+1}=x_k-H_k^{-1}\nabla f(x_k)$：
\[
  e_{k+1}=e_k-H_k^{-1}(He_k+R_k)
  =\bigl(I-H_k^{-1}H\bigr)e_k-H_k^{-1}R_k
  =H_k^{-1}E_ke_k-H_k^{-1}R_k .
\]
由 $\|E_k\|\le M\|e_k\|\le\mu/2$ 与 Weyl 不等式，
$\lambda_{\min}(H_k)\ge\mu-\|E_k\|\ge\mu/2$，故 $\|H_k^{-1}\|\le2/\mu$。
于是
\[
  \|e_{k+1}\|\le\frac2\mu\Bigl(M\|e_k\|\cdot\|e_k\|+\frac M2\|e_k\|^2\Bigr)
  =\frac{3M}{\mu}\|e_k\|^2 . \qedhere
\]
\end{proof}

\begin{corollary}[三个前提缺一不可，缺哪一条就退化成什么]
\label{cor:newton-degrade}
\begin{enumerate}
  \item \textbf{离开解附近}：定理 \ref{thm:newton-quadratic} 的常数项 $\frac{3M}\mu\|e_k\|$
        本身是 $O(\|e_k\|)$，远一步就不成立。实测四次型上
        $e_0=1.414$、$q_0=e_1/e_0^2=0.1076$，而渐近常数是
        $C\approx9.905\times10^{-5}$，$q_0/C=1087$：
        第一步比渐近常数大三个数量级。
  \item \textbf{不整步}：若每步只允许走 $\alpha_{\text{cap}}<1$ 倍牛顿步，
        解附近 $H_k\to H$，于是 $e_{k+1}=(1-\alpha_{\text{cap}})e_k+O(\|e_k\|^2)$，
        退化为线性收敛，因子 $1-\alpha_{\text{cap}}$。
        实测 $\alpha_{\text{cap}}=0.75/0.5/0.25$ 对应线性因子
        $0.25/0.5/0.7506$（与 $1-\alpha_{\text{cap}}$ 的偏差小于 $0.05$）。
  \item \textbf{Hessian 不正定}：见下面的定理 \ref{thm:newton-saddle}：
        牛顿法照样一步到达驻点，只是那个驻点是鞍点。
\end{enumerate}
\end{corollary}

\begin{proof}
(1) 定理 \ref{thm:newton-quadratic} 给出的收缩因子是
$\frac{3M}{\mu}\|e_k\|$。只有当
$\|e_k\|<\frac{\mu}{3M}$ 时它才小于 $1$，此后每步误差才被平方级压小；
在此之前 $\|e_{k+1}\|$ 甚至可能暂时变大（$e_0=1.414$ 时该因子是 $1087C/\|e_0\|$
量级，实测 $q_0/C=1087$）。

(2) 设 $x_{k+1}=x_k-\alpha_{\text{cap}}H_k^{-1}\nabla f(x_k)$，$H_k=\nabla^2f(x_k)$。
沿用 (1) 的记号 $\nabla f(x_k)=He_k+R_k$（$\|R_k\|\le\frac M2\|e_k\|^2$）、
$H_k=H+E_k$（$\|E_k\|\le M\|e_k\|$），则
\[
  e_{k+1}=(1-\alpha_{\text{cap}})e_k
  +\alpha_{\text{cap}}\bigl(H_k^{-1}E_k e_k-H_k^{-1}R_k\bigr),
\]
其中最后一项是 $O(\|e_k\|^2)$。于是解附近
$\|e_{k+1}\|=(1-\alpha_{\text{cap}})\|e_k\|+O(\|e_k\|^2)$，
即线性收敛，因子 $1-\alpha_{\text{cap}}$（实测 $0.25/0.5/0.7506$，
与 $1-\alpha_{\text{cap}}$ 的偏差小于 $0.05$）。

(3) 由定理 \ref{thm:newton-saddle}，$\lambda_{\min}<0$ 时 $x^\star$ 是鞍点，
牛顿法照样一步到达；"收敛"与"收敛到极小点"是两件事。
\end{proof}

\begin{theorem}[牛顿法不检查曲率符号：鞍点]
\label{thm:newton-saddle}
设 $f(x)=\tfrac12x^{\mathsf T}Hx$，$H$ 非奇异且 $\lambda_{\min}(H)<0$，
对应单位特征向量 $u$。则：
\begin{enumerate}
  \item 整步牛顿法从任意 $x_0$ 出发一步到达 $x^\star=0$；
  \item 定步长梯度下降在特征方向上逐个递推，
        $u^{\mathsf T}e_{k+1}=(1-\alpha\lambda_{\min})\,u^{\mathsf T}e_k$，
        故负曲率分量以 $|1-\alpha\lambda_{\min}|=1+|\lambda_{\min}|\alpha>1$ 增长；
  \item 要在 $N$ 步内离开半径 $R$ 的邻域，扰动沿 $\pm u$ 的分量必须满足
        \[
          \delta^\star=R\,(1+|\lambda_{\min}|\alpha)^{-N};
        \]
  \item 若扰动与 $u$ 正交，则 $u^{\mathsf T}e_k\equiv0$，永远不逃逸。
\end{enumerate}
\end{theorem}

\begin{proof}
(1) 与定理 \ref{thm:newton-one-step} 同，取 $A=H$、$b=0$。

(2) 由定理 \ref{thm:gd-recurrence}，$e_{k+1}=(I-\alpha H)e_k$，
沿 $u$ 投影得 $u^{\mathsf T}e_{k+1}=(1-\alpha\lambda_{\min})u^{\mathsf T}e_k$；
$\alpha>0$、$\lambda_{\min}<0$ 时 $1-\alpha\lambda_{\min}>1$。

(3) $u^{\mathsf T}e_k=(1-\alpha\lambda_{\min})^k\,u^{\mathsf T}e_0$，
要求存在 $k\le N$ 使 $\|e_k\|\ge R$；在纯 $u$ 方向扰动下
$\|e_k\|=|u^{\mathsf T}e_k|$，故最小初始分量是 $R(1+|\lambda_{\min}|\alpha)^{-N}$。

(4) 若 $u^{\mathsf T}e_0=0$，由 (2) 的递推立得 $u^{\mathsf T}e_k=0$ 对所有 $k$ 成立；
其余方向按 $|1-\alpha\lambda_i|<1$ 收缩，迭代点收敛到鞍点。
\end{proof}

实测（\texttt{test\_escape\_threshold\_matches\_analytic\_formula}）：$\lambda_{\min}=-0.05$、
$\alpha=1$、$R=1$、$N=120$ 时二分出来的阈值是 $2.86605\times10^{-3}$，
解析值 $(1.05)^{-120}=2.86605\times10^{-3}$，相对差 $5.3\times10^{-10}$；
沿 $45^\circ$ 方向（与 $u_{\min}$ 的投影为 $\tfrac1{\sqrt2}$）阈值放大
$\sqrt2$ 倍（相对差小于 $10^{-5}$）。
与 $u_{\min}$ 正交的方向上，把试探幅度放大到 $10^{6}$ 倍仍然测不到逃逸
（\texttt{test\_orthogonal\_perturbation\_never\_escapes}）——
这时阈值必须写成"没有测到"，不是 $0$。

\section{线搜索：步长是怎么被定下来的}

牛顿方向只给出"往哪走"，走多远由线搜索决定。下面这条定理是 Armijo 回溯的全部理论基础，
它也解释了为什么回溯**一定**会在有限步内停下。

\begin{theorem}[Armijo 回溯的终止性与步长下界]
\label{thm:armijo}
设 $f\in C^1$，$\nabla f$ 在 $x$ 到 $x+\alpha_0p$ 的线段上 Lipschitz 连续、常数 $L$，
$g=\nabla f(x)$，$p$ 是下降方向（$g^{\mathsf T}p<0$）。
取 $c_1\in(0,1)$、$\rho\in(0,1)$，令 $\alpha=\alpha_0\rho^m$（$m=0,1,2,\dots$）为第一个满足
\[
  f(x+\alpha p)\le f(x)+c_1\alpha\,g^{\mathsf T}p
\]
的步长。则这样的 $m$ 一定存在，且
\[
  \alpha\ \ge\ \min\left\{\alpha_0,\ \frac{2\rho(1-c_1)\,|g^{\mathsf T}p|}{L\|p\|^2}\right\}.
\]
\end{theorem}

\begin{proof}
记 $\phi(t)=f(x+tp)$，则 $\phi'(t)=\nabla f(x+tp)^{\mathsf T}p$。由 Lipschitz 连续性，
\[
  \phi(\alpha)-\phi(0)-\alpha\phi'(0)=\int_0^\alpha\bigl(\phi'(t)-\phi'(0)\bigr)\,dt
  \le\int_0^\alpha Lt\|p\|^2\,dt=\frac L2\alpha^2\|p\|^2 ,
\]
即（下降引理）$f(x+\alpha p)\le f(x)+\alpha g^{\mathsf T}p+\frac L2\alpha^2\|p\|^2$。
于是只要
\[
  \frac L2\alpha^2\|p\|^2\le -(1-c_1)\alpha\,g^{\mathsf T}p
  \iff \alpha\le\bar\alpha:=\frac{2(1-c_1)|g^{\mathsf T}p|}{L\|p\|^2},
\]
Armijo 条件就成立。$\bar\alpha>0$，而 $\alpha_0\rho^m\to0$，故回溯必然在有限步内
（至多 $\lceil\log_\rho(\bar\alpha/\alpha_0)\rceil+1$ 步）达到 $\alpha\le\bar\alpha$ 并接受。
若 $\alpha_0\le\bar\alpha$ 则第一步就接受，$\alpha=\alpha_0$；
否则被接受的那一步的前一步 $\alpha/\rho$ 必然 $>\bar\alpha$（否则会更早接受），
于是 $\alpha>\rho\bar\alpha$。两种情形合起来即为所证下界。
\end{proof}

\begin{remark}[强 Wolfe 曲率条件的数值边界]
强 Wolfe 的第二个条件 $|\phi'(\alpha)|\le c_2|\phi'(0)|$ 要求把方向导数分辨到
$c_2|g^{\mathsf T}p|$ 的水平。而浮点算出的 $\phi'(\alpha)=g(x+\alpha p)^{\mathsf T}p$
带有量级 $u\,L\|x\|\,\|p\|$ 的绝对误差：当 $c_2|g^{\mathsf T}p|$ 落到这个水平以下时，
曲率条件可能对**任何** $\alpha$ 都不成立（这不是算法错，是判据在该精度下没有意义）。
所以本章的 \texttt{wolfe\_line\_search} 把两件事分开返回：
\texttt{ok}（Armijo 成立、这一步可用）与 \texttt{wolfe}（曲率条件是否也满足），
并把后者如实记进 \texttt{wolfe\_ok} 数组，绝不把"曲率条件达不到"当成"线搜索失败"而中止迭代。
实测：本章的算例上曲率条件在被接受的每一步都成立
（Rosenbrock 36 步、4 维 $\kappa=1000$ 二次型 9 步、四次型 20 步，\texttt{wolfe\_ok} 全为 $1$）。
\end{remark}

\begin{remark}[BFGS 前几步的步长]
实测 4 维 $\kappa=1000$ 二次型上 BFGS 的前三步步长是
$9.77\times10^{-4},\ 1.56\times10^{-2},\ 0.125$，之后立刻变成 $1.0$：
一开始曲率信息不足，线搜索只能给出很小的步长；三次更新之后就接近牛顿步。
这也是"拟牛顿在前几步很慢、之后很快"的定量版本。
\end{remark}

\section{拟牛顿：割线方程与正定性}

\begin{theorem}[BFGS 更新精确满足最新一条割线方程]
\label{thm:bfgs-secant}
设 $B_k$ 对称、$s_k=x_{k+1}-x_k$、$y_k=\nabla f(x_{k+1})-\nabla f(x_k)$、
$y_k^{\mathsf T}s_k\ne0$，且
\[
  B_{k+1}=B_k-\frac{B_ks_ks_k^{\mathsf T}B_k}{s_k^{\mathsf T}B_ks_k}
  +\frac{y_ky_k^{\mathsf T}}{y_k^{\mathsf T}s_k}.
\]
则 $B_{k+1}s_k=y_k$。
\end{theorem}

\begin{proof}
逐项右乘 $s_k$：
\[
  B_{k+1}s_k=B_ks_k-\frac{B_ks_k\,(s_k^{\mathsf T}B_ks_k)}{s_k^{\mathsf T}B_ks_k}
  +\frac{y_k\,(y_k^{\mathsf T}s_k)}{y_k^{\mathsf T}s_k}
  =B_ks_k-B_ks_k+y_k=y_k . \qedhere
\]
\end{proof}

\begin{theorem}[$s^{\mathsf T}y>0$ 时 BFGS 保持正定]
\label{thm:bfgs-pd}
若 $B_k\succ0$ 且 $s_k^{\mathsf T}y_k>0$，则定理 \ref{thm:bfgs-secant} 的 $B_{k+1}\succ0$。
\end{theorem}

\begin{proof}
对任意 $v\ne0$，记 $a=s_k^{\mathsf T}B_ks_k>0$（因 $B_k\succ0$、$s_k\ne0$）：
\[
  v^{\mathsf T}B_{k+1}v
  =\underbrace{\frac{(v^{\mathsf T}B_kv)\,a-(v^{\mathsf T}B_ks_k)^2}{a}}_{\text{(I)}}
  +\underbrace{\frac{(v^{\mathsf T}y_k)^2}{s_k^{\mathsf T}y_k}}_{\text{(II)}} .
\]
(I) $\ge0$ 由 $B_k$ 内积下的 Cauchy-Schwarz 不等式；(II) $\ge0$ 由 $s_k^{\mathsf T}y_k>0$。
若 $v^{\mathsf T}B_{k+1}v=0$，则 (II) $=0$ 给出 $v^{\mathsf T}y_k=0$；(I) $=0$ 要求
Cauchy-Schwarz 取等，即 $v$ 与 $s_k$ 共线，$v=ts_k$（$t\ne0$），
于是 $v^{\mathsf T}y_k=ts_k^{\mathsf T}y_k\ne0$，矛盾。故 $v^{\mathsf T}B_{k+1}v>0$。
\end{proof}

\begin{theorem}[L-BFGS 的两循环递推与"只有最新一对精确"]
\label{thm:lbfgs}
记 $\rho_i=1/(y_i^{\mathsf T}s_i)$、$V_i=I-\rho_iy_is_i^{\mathsf T}$，
$H_k$ 由 $H_0=\gamma I$ 与 $H_{i+1}=V_i^{\mathsf T}H_iV_i+\rho_is_is_i^{\mathsf T}$
递推得到（只保留最近 $m$ 对）。则：
\begin{enumerate}
  \item 两循环算法 $r\gets H_kg$ 逐步计算的就是这个 $H_kg$；
  \item $H_ky_k=s_k$（最新一对精确满足）；
  \item 更老的配对一般不满足：$H_ky_i\ne s_i$（$i<k$）。
\end{enumerate}
\end{theorem}

\begin{proof}
(1) 把 $H_k$ 的递推展开成
$H_k=\sum_{i}\rho_i\bigl(\prod_{j>i}V_j^{\mathsf T}\bigr)s_is_i^{\mathsf T}
\bigl(\prod_{j>i}V_j\bigr)+\bigl(\prod_jV_j^{\mathsf T}\bigr)H_0
\bigl(\prod_jV_j\bigr)$，
从最内层往外乘即得两循环：第一圈从最新往最老算
$\alpha_i=\rho_is_i^{\mathsf T}r$、$r\gets r-\alpha_iy_i$，
中间乘 $H_0=\gamma I$，第二圈从最老往最新把 $\rho_is_i\alpha_i$ 加回去。

(2) 取 $r=y_k$。第一圈的第一个（最新）$\alpha_k=\rho_ks_k^{\mathsf T}y_k=1$，
随后 $r\gets y_k-y_k=0$，故其余 $\alpha_i=\rho_is_i^{\mathsf T}0=0$；
中间 $r\gets\gamma r=0$；第二圈只有 $i=k$ 的项 $\rho_ks_k\alpha_k=s_k$ 非零，
其余项为 $0$。所以 $H_ky_k=s_k$。

(3) $H_k$ 只由 $(s_j,y_j)_{j\le k}$ 决定，且由 (2) 在更早时刻
$H_{k-1}y_{k-1}=s_{k-1}$；但 $H_k$ 相对 $H_{k-1}$ 又做了一次更新，
对 $y_{k-1}$ 的作用不再等于 $s_{k-1}$。实测量级见下。
\end{proof}

实测（\texttt{test\_lbfgs\_secant\_residual\_grows\_with\_pair\_age}）：$m=2$、$n=4$、
$\kappa=1000$ 时按年龄（新到老）的逆割线残差是 $0,\ 0.234$；
$m=5$ 时是 $0,\ 1.688,\ 0.384,\ 0.079,\ 0.132$；
$m=20\ge n$ 时最新一对是 $2.8\times10^{-16}$、最老一对也降到 $0.090$。
注意定理 \ref{thm:bfgs-secant} 说的是"更新后满足、更新前不满足"，
而这里说的是"最新一对满足、更老的只近似"——两个"只有"是不同的两件事，
混在一起说就会得出"L-BFGS 不满足割线方程"这种错误结论。

\section{信赖域}

\begin{theorem}[精确模型下 $\rho\equiv1$]
\label{thm:tr-rho-one}
若模型 $m_k$ 与被优化的 $f$ 在信赖域内处处相等（例如目标是二次函数、
模型取二阶 Taylor 且二阶项就是精确 Hessian），则对任意试探步 $p$ 有
$\rho_k=1$。于是：
\begin{enumerate}
  \item 没有任何一步会被拒绝（拒绝需要 $\rho_k<\eta=0.15$）;
  \item 半径永不收缩（收缩需要 $\rho_k<t_1=0.25$），$\Delta_k$ 单调不减。
\end{enumerate}
\end{theorem}

\begin{proof}
$\rho_k=\dfrac{f(x_k)-f(x_k+p)}{m_k(0)-m_k(p)}=\dfrac{m_k(0)-m_k(p)}{m_k(0)-m_k(p)}=1$
（分母非零时）。$1>\eta$ 且 $1>t_1$，故既不拒绝也不收缩；
扩张只在 $\|p\|=\Delta_k$ 且 $\rho_k>t_2$ 时发生，因此 $\Delta_k$ 不会变小。
\end{proof}

实测：$\kappa=100$ 的二次型上 $\rho$ 序列是 $(1.000000,1.000000)$，
半径 $1\to2\to2$，拒绝 $0$ 次、负曲率 $0$ 次。
所以\textbf{在二次型上看不到任何"自适应"}——$\rho\equiv1$ 意味着模型完全准，
没有东西需要适应。要演示自适应必须换非二次目标：Rosenbrock 从 $(-1.2,1)$ 出发，
$22$ 步里 $3$ 次拒绝、$\rho$ 出现 $-4.25$、半径在 $0.125$ 与 $1.0$ 之间来回；
从 $(0.5,3)$ 出发还会触发 $3$ 次负曲率分支（该点 Hessian 行列式为负）。

\begin{theorem}[Cauchy 点的模型下降下界]
\label{thm:cauchy}
设 $B\succ0$、$g=\nabla f(x_k)\ne0$，Cauchy 点指沿 $-g$ 方向使模型
$m(p)=g^{\mathsf T}p+\tfrac12p^{\mathsf T}Bp$ 最小的点（限制在 $\|p\|\le\Delta$ 内）。
则其模型下降量满足
\[
  m(0)-m(p_C)\ \ge\ \tfrac12\|g\|\min\!\left(\Delta,\ \frac{\|g\|}{\|B\|}\right).
\]
\end{theorem}

\begin{proof}
沿 $-g$ 单位方向写 $p=-t\,g/\|g\|$（$0\le t\le\Delta$），则
\[
  m(p)=-t\|g\|+\frac{t^2}{2}\,\frac{g^{\mathsf T}Bg}{\|g\|^2}.
\]
取 $t^\ast=\min\bigl(\Delta,\ \|g\|/\|B\|\bigr)$。
\emph{情形一} $\Delta\le\|g\|/\|B\|$：取 $t=\Delta$，
$m(0)-m(p_C)\ge\|g\|\Delta-\frac12\Delta^2\|B\|
=\|g\|\Delta\bigl(1-\frac{\Delta\|B\|}{2\|g\|}\bigr)
\ge\frac12\|g\|\Delta$，
最后一步用了 $\Delta\|B\|\le\|g\|$。
\emph{情形二} $\Delta>\|g\|/\|B\|$：取 $t=\|g\|/\|B\|\le\Delta$，
$m(0)-m(p_C)\ge\frac{\|g\|^2}{\|B\|}-\frac12\frac{\|g\|^2}{\|B\|}
=\frac12\frac{\|g\|^2}{\|B\|}$。
两种情形合起来就是
$m(0)-m(p_C)\ge\frac12\|g\|\min(\Delta,\|g\|/\|B\|)$。
\end{proof}

\begin{remark}[这个下界是干什么用的]
它保证"只要梯度不为零，模型至少能下降 $\frac12\|g\|\min(\Delta,\|g\|/\|B\|)>0$"，
这是信赖域方法全局收敛性证明里唯一需要的定量输入。
本章的可视化用的是 Steihaug-CG 解（未必是 Cauchy 点），
但每步的 $\rho_k$ 与"被拒绝的步"都记录在帧里，
\texttt{meta.n\_rejected}、\texttt{meta.ratio\_min} 就是这条下界在实际算法里的影子。
\end{remark}

\section{非线性共轭梯度}

\begin{theorem}[精确线搜索下 CG 的有限步终止与共轭性]
\label{thm:cg}
设 $A\succ0$、$f$ 为二次型，CG 取 $p_0=-g_0$、
$p_{k+1}=-g_{k+1}+\beta_kp_k$、$\beta_k=\dfrac{g_{k+1}^{\mathsf T}g_{k+1}}{g_k^{\mathsf T}g_k}$
（Fletcher-Reeves），步长由精确线搜索 $\alpha_k=\arg\min_t f(x_k+tp_k)$ 给出。则
\[
  g_i^{\mathsf T}g_j=0\ (i\ne j),\qquad p_i^{\mathsf T}Ap_j=0\ (i\ne j),
\]
且 $x_{k+1}$ 是 $f$ 在仿射子空间 $x_0+\mathrm{span}\{g_0,\dots,g_k\}$ 上的极小点；
特别地 $x_n=x^\star$。
\end{theorem}

\begin{proof}
对 $k$ 归纳，同时维护三条：
(i) $g_i^{\mathsf T}g_j=0$（$0\le i<j\le k$）；
(ii) $p_i^{\mathsf T}Ap_j=0$（$0\le i<j\le k$）；
(iii) $\mathrm{span}\{p_0,\dots,p_k\}=\mathrm{span}\{g_0,\dots,g_k\}
=\mathcal K_{k+1}(A,g_0)$。
$k=0$ 时三条都平凡。归纳步用两个恒等式：精确线搜索给出 $p_k^{\mathsf T}g_{k+1}=0$，
配合 $g_{k+1}=g_k+\alpha_kAp_k$ 与 $p_k^{\mathsf T}g_k=-\|g_k\|^2$ 得到
\[
  \alpha_k=\frac{\|g_k\|^2}{p_k^{\mathsf T}Ap_k};
\]
以及由 (iii) $g_k=-p_k+\beta_{k-1}p_{k-1}$。

\emph{正交性}：对 $i\le k$ 写 $g_{k+1}=g_i+\sum_{j=i}^{k}\alpha_jAp_j$，则
$g_{k+1}^{\mathsf T}g_i=\|g_i\|^2+\sum_{j=i}^{k}\alpha_j\,p_j^{\mathsf T}Ag_i$。
由 (iii)，$g_i\in\mathrm{span}\{p_0,\dots,p_i\}$，故
$Ag_i\in\mathrm{span}\{Ap_0,\dots,Ap_i\}$，而 (ii) 说 $p_j$ 与 $Ap_m$
在 $j>m$ 时正交，于是和式里 $j>i$ 的项全为零，只剩
$\|g_i\|^2+\alpha_ip_i^{\mathsf T}Ag_i$；
再用 $g_i=-p_i+\beta_{i-1}p_{i-1}$ 与 (ii) 得 $p_i^{\mathsf T}Ag_i=-p_i^{\mathsf T}Ap_i$，
代入 $\alpha_i$ 的表达式即得 $g_{k+1}^{\mathsf T}g_i=0$。

\emph{共轭性}：对 $j\le k$，
$p_{k+1}^{\mathsf T}Ap_j=-g_{k+1}^{\mathsf T}Ap_j+\beta_kp_k^{\mathsf T}Ap_j$。
$j<k$ 时第二项由 (ii) 为零，第一项用 $Ap_j=(g_{j+1}-g_j)/\alpha_j$
与刚证的正交性为零。$j=k$ 时两项分别是
$g_{k+1}^{\mathsf T}Ap_k=\|g_{k+1}\|^2/\alpha_k$ 与
$\beta_kp_k^{\mathsf T}Ap_k=\beta_k\|g_k\|^2/\alpha_k$，
而 $\beta_k=\|g_{k+1}\|^2/\|g_k\|^2$，故和为零。

(iii) 由 $p_{k+1}=-g_{k+1}+\beta_kp_k$ 与 $g_{k+1}=g_k+\alpha_kAp_k$ 直接得到。
最后：在 $A$-共轭方向组上依次做精确一维极小，得到的点就是整个
$\mathrm{span}\{p_0,\dots,p_k\}$ 上的极小点（这是二次函数的标准性质：
$A$-共轭方向组是该子空间的一组正交基，各方向的极小互不干扰）。
于是 $x_n$ 是 $f$ 在整个 $\mathbb R^n$ 上的极小点，
即 $\nabla f(x_n)=g_n=0$，而 (i) 说 $\{g_i\}$ 两两正交、
最多只有 $n$ 个非零元，所以 $g_n=0$ 必在 $n$ 步之内发生。
\end{proof}

实测（\texttt{test\_cg\_exact\_line\_search\_terminates\_in\_n\_steps}）：
$(n,\kappa)=(4,10^3),(6,10^2),(8,10)$ 上都恰好 $n$ 步，
归一化共轭偏差 $\max_{i\ne j}|p_i^{\mathsf T}Ap_j|/(\|p_i\|_A\|p_j\|_A)\le1.7\times10^{-10}$。
但定理 \ref{thm:cg} 的两个前提都是"精确"的：
换成 Wolfe 不精确线搜索后 $n=4$ 要 $72$ 步；
$\kappa=10^5$、$n=6$ 时\textbf{即使精确线搜索}也要 $8$ 步，
且共轭偏差高达 $0.995$——共轭性已经被浮点误差摧毁。

\section{四个可视化怎么读}

\begin{enumerate}
  \item \texttt{descent\_race}：同一张等高线上四条路径。
        \texttt{step\_scale} 是 $\alpha/\alpha^\star$，
        它的临界值 $1.0$ 是"最优"，而发散线在 $1+1/\kappa$（$\kappa=100$ 时是 $1.01$）——
        两个临界值挨得这么近，正是这一页要让人看见的事。
        副轴的误差曲线在 $10^{-17}$ 处截平（\texttt{meta.n\_exact\_end} 给出真正到 $0$ 的点数），
        发散的路径用尾部重复末点补位并在 \texttt{strings} 里说明。
  \item \texttt{convergence\_rates}：横轴是迭代数，纵轴是对数误差。
        要看的是 $q_k=e_{k+1}/e_k^2$ 这条曲线**平不平**：
        默认参数下有 2 步平（$C\approx9.9\times10^{-5}$），
        然后突然冲到 $10^{10}$ 以上——那是舍入地板（\texttt{meta.newton\_floor\_start}），
        不是"收敛越来越快"。把 \texttt{alpha\_cap} 拖到 $1$ 以下，平地当场消失。
  \item \texttt{saddle\_escape}：鞍点上的等高线是双曲线族。
        \texttt{direction} 选 \texttt{orthogonal} 时轨迹会滑回鞍点；
        选 \texttt{neg\_curv} 时按 $(1+|\lambda_-|\alpha)^k$ 长出去。
        \texttt{meta.threshold\_neg\_curv} 与 \texttt{meta.threshold\_random} 的比值
        就是 $1/|\langle u_{\min},d\rangle|$。
  \item \texttt{trust\_region\_radius}：半径 $\Delta_k$、比值 $\rho_k$、
        接受掩码、负曲率掩码四条曲线一起看。
        起点选 \texttt{indefinite} 才会看到 Steihaug-CG 走负曲率分支。
        $\rho_k$ 可以为负（这一步会让 $f$ 上升，被 $\eta$ 挡住），
        \texttt{ratio\_history} 里的 $0$ 是占位符，必须配合 \texttt{ratio\_valid} 读。
\end{enumerate}

% 本文件由 tools/build_book.py 从 dist/reports/assertions.json 生成，请勿手改。
\begin{longtable}{@{}p{0.06\textwidth}p{0.32\textwidth}p{0.22\textwidth}p{0.16\textwidth}@{}}
\caption{章节 5.1 数值断言（生成时间 2026-10-19T16:08:50-0400）}\\
\toprule
类别 & 断言 & 理论值 & 实测值 \\
\midrule
\endfirsthead
\toprule
类别 & 断言 & 理论值 & 实测值 \\
\midrule
\endhead
\bottomrule
\endlastfoot

A & 端点组合帧数与 JSON 里的 NaN/Infinity 计数 & NaN/\allowbreak{}Infinity 计数 =\allowbreak{} 0 & 69 帧，计数 0(\allowbreak{}descent\_race: 16 个数组、\allowbreak{}20 个 meta 字段) \checkmark \\
A & 对数轴曲线最小值 > 0、掩码形状与 dtype、strings 里的地板说明 & 全部满足 & descent\_race.error\_gd min=\allowbreak{}4.3e-01；descent\_race.error\_newton min=\allowbreak{}1.3e-14；convergence\_rates.error\_gd min=\allowbreak{}1.0e+00；convergence\_rates.error\_newton min=\allowbreak{}1.0e-17；saddle\_escape.residual min=\allowbreak{}4.0e-04；saddle\_escape.error\_to\_saddle min=\allowbreak{}2.0e-02；q\_newton\_valid dtype=\allowbreak{}bool \checkmark \\
A & 发散帧的三个诚实字段 & diverged\_gd=\allowbreak{}True(\allowbreak{}放大 1/\allowbreak{}u 倍即判发散)、\allowbreak{}overflow=\allowbreak{}True、\allowbreak{}stop\_gd=\allowbreak{}'diverged' & diverged\_gd=\allowbreak{}True，overflow=\allowbreak{}True，stop\_gd=\allowbreak{}'diverged'，$\alpha$/\allowbreak{}$\alpha$\_stable=\allowbreak{}1.5998 \checkmark \\
A & 发散帧里 error\_\allowbreak{}gd 的长度与有限性 & 长度 =\allowbreak{} max\_iter+1 =\allowbreak{} 401，全有限(\allowbreak{}尾部补位) & 长度 401，max=\allowbreak{}3.027e+15，尾部 3 个 [3.027e+15 3.027e+15 3.027e+15] \checkmark \\
A & meta.n\_\allowbreak{}iter\_\allowbreak{}gd 与核心 gd(...) 的迭代数 & 相等 & 帧 60 vs 核心 60(\allowbreak{}$\alpha$=\allowbreak{}1.980198e-02) \checkmark \\
A & meta.kappa\_\allowbreak{}measured 与实测谱 & =\allowbreak{} $\kappa$(\allowbreak{}Jacobi 实测，相对 1e-12) & 100.0000000000(\allowbreak{}给定 100.0) \checkmark \\
A & 帧 meta.newton\_\allowbreak{}quadratic\_\allowbreak{}constant 与核心实测 C & 相等(\allowbreak{}同一窗口定义 + 同一初值/\allowbreak{}容差) & 帧 9.9047694185682e-05 vs 核心 9.9047694185682e-05(\allowbreak{}enter\_iter=\allowbreak{}2) \checkmark \\
A & $\beta$=0(退化成二次型)时的常数与说明 & 常数 =\allowbreak{} None，且 strings 里有解释 & 常数 None，enter\_iter None，string 提到 null=\allowbreak{}True \checkmark \\
A & alpha\_\allowbreak{}cap=1 时的尾部阶 / 窗口 / 地板起点 / q 最大值 & 阶 $\ge$ 1.8、\allowbreak{}窗口 $\ge$ 2、\allowbreak{}常数有限、\allowbreak{}地板起点有限、\allowbreak{}q 冲到 $\ge$ 1e50 & 阶 2.999，窗口 2，C=\allowbreak{}9.9047694185682e-05，地板起点 k=\allowbreak{}4，q\_max=\allowbreak{}3.35e+122 \checkmark \\
A & alpha\_\allowbreak{}cap=0.5 时的尾部阶 / 窗口 / 常数 & 阶 $\approx$ 1(\allowbreak{}线性)、\allowbreak{}窗口 0、\allowbreak{}常数 None & 阶 1.000，窗口 0，C=\allowbreak{}None，步数 30 \checkmark \\
A & 两个方向上的扰动末距($\delta$=0.1, 200 步) & neg\_curv 远大于 1 远大于 orthogonal(\allowbreak{}差 $\ge$ 1e3 倍) & neg\_curv 1.729e+03 vs orthogonal 0.000e+00(\allowbreak{}比 1.73e+303) \checkmark \\
A & 实测 $\lambda$\_\allowbreak{}min 与解析阈值的对照 & $\lambda$\_min =\allowbreak{} -0.05，阈值 $\approx$ 2.87e-3(\allowbreak{}N=\allowbreak{}96 时的解析值) & $\lambda$\_min=\allowbreak{}-0.050000，u\_min=\allowbreak{}(\allowbreak{}1.000,\allowbreak{}0.000)，阈值 5.7830e-05，解析 5.7828e-05 \checkmark \\
A & contour\_\allowbreak{}z 与逐点求值的最大相对差 & $\le$ 1e-14 & 0.000e+00(\allowbreak{}网格 (\allowbreak{}61,\allowbreak{} 61)) \checkmark \\
A & 不定起点上的负曲率分支 / 拒绝次数 / 负 $\rho$ & n\_negative\_curv $\ge$ 1、\allowbreak{}n\_rejected $\ge$ 1、\allowbreak{}min $\rho$ < 0 & 负曲率 3 次，拒绝 3 次，min $\rho$=\allowbreak{}-1.801，接受掩码里 0 的个数 9 \checkmark \\
A & 半径的收缩与扩张范围 & radius\_min < $\Delta$0=\allowbreak{}1.0 < radius\_max & [0.2500,\allowbreak{} 2.0000]，终值 1.0000 \checkmark \\
A & 轨迹 / 半径历史 / 梯度范数与核心的最大偏差 & 0(\allowbreak{}同一份数组) & 轨迹 0.0e+00，半径 0.0e+00，梯度 0.0e+00(\allowbreak{}21 步，帧长 25) \checkmark \\
A & meta.n\_\allowbreak{}accepted + meta.n\_\allowbreak{}rejected 与总步数 & 和 =\allowbreak{} trace 的步数(\allowbreak{}被拒绝的步不画进轨迹，但计数里在) & 21 + 3 =\allowbreak{} 24(\allowbreak{}max\_iter=\allowbreak{}24，实际 21 步) \checkmark \\
A & 进入二次区的步号 / 窗口内 q 的散布 & 窗口 $\ge$ 2 点、\allowbreak{}散布 $\le$ 2.5(\allowbreak{}相邻比值定义)，步号有限 & enter\_iter=\allowbreak{}2，窗口 [9.905e-05 6.835e-05]，散布 1.449，C$\approx$9.905e-05 \checkmark \\
A & 第一步 e1/e0\textasciicircum{}2 $\div$ 渐近常数 C & $\ge$ 100(\allowbreak{}远处不是二次收敛) & 1.087e+03(\allowbreak{}第一步 1.076e-01) \checkmark \\
A & 渐近二次常数 C 与牛顿步数 & C 有限且 > 0；步数 $\le$ 8 & C=\allowbreak{}9.904769e-05，4 步到 6.454e-25 \checkmark \\
A & 线性因子 - (1-$\alpha$\_\allowbreak{}cap) 的最大偏差 & $\le$ 0.05 & 0.0000(\allowbreak{}$\alpha$\_cap=\allowbreak{}0.75: 实测 0.2500 /\allowbreak{} 理论 0.2500(\allowbreak{}25 步)；$\alpha$\_cap=\allowbreak{}0.5: 实测 0.5000 /\allowbreak{} 理论 0.5000(\allowbreak{}50 步)；$\alpha$\_cap=\allowbreak{}0.25: 实测 0.7500 /\allowbreak{} 理论 0.7500(\allowbreak{}121 步)) \checkmark \\
A & $\alpha$\_\allowbreak{}cap=1 与 $\alpha$\_\allowbreak{}cap=0.5 的步数比 & $\ge$ 5(\allowbreak{}二次 vs 线性) & 4 : 50(\allowbreak{}比 12.5) \checkmark \\
A & 正定二次型上牛顿一步之后的 $\|$x1 - x*$\|$ & $\le$ $\kappa$$\cdot$u(\allowbreak{}一次求解的舍入量级) & 1.275e-14(\allowbreak{}$\kappa$=\allowbreak{}100.0，一步，正定标志 [1.]) \checkmark \\
A & 鞍点上牛顿一步之后的 dist 与 Hessian 正定标志 & dist $\approx$ 0(\allowbreak{}到驻点)，但 $\lambda$\_min<0、\allowbreak{}正定标志 0 & dist=\allowbreak{}0.000e+00，$\lambda$\_min=\allowbreak{}-0.500，正定标志 [0.] \checkmark \\
A & BFGS 更新后割线残差 max $\|$B s - y$\|$/$\|$y$\|$ & $\le$ 100u $\approx$ 1.1e-14 & 1.055e-14(\allowbreak{}4 维 $\kappa$=\allowbreak{}1000 二次型: 更新后 7.49e-15 /\allowbreak{} 更新前 0.999(\allowbreak{}9 步)；四次型: 更新后 1.05e-14 /\allowbreak{} 更新前 0.990(\allowbreak{}7 步)) \checkmark \\
A & BFGS 更新前割线残差 max(同一批步) & $\ge$ 0.5(\allowbreak{}更新不是多余的) & 0.999 \checkmark \\
A & L-BFGS(m=2) 最新一对的逆割线残差 & =\allowbreak{} O(\allowbreak{}u)(\allowbreak{}构造上精确) & 0.000e+00 \checkmark \\
A & L-BFGS(m=2) 更老一对的逆割线残差 & 远大于 最新一对(\allowbreak{}记忆没覆盖到) & 0.234(\allowbreak{}比值 2.339e+299) \checkmark \\
A & CG(精确线搜索) 迭代数 - n($\kappa$ $\le$ 1e3) & 0(\allowbreak{}有限步终止：至多 n 步精确到达) & n=\allowbreak{}4,\allowbreak{}$\kappa$=\allowbreak{}1000: 4 步，共轭偏差 1.71e-10；n=\allowbreak{}6,\allowbreak{}$\kappa$=\allowbreak{}100: 6 步，共轭偏差 3.56e-10；n=\allowbreak{}8,\allowbreak{}$\kappa$=\allowbreak{}10: 8 步，共轭偏差 5.15e-12 \checkmark \\
A & 归一化 A-共轭偏差 max |p\_\allowbreak{}i\textasciicircum{}TAp\_\allowbreak{}j|/($\|$p\_\allowbreak{}i$\|$\_\allowbreak{}A$\|$p\_\allowbreak{}j$\|$\_\allowbreak{}A) & $\le$ 1e-6(\allowbreak{}$\kappa$ $\le$ 1e3 时舍入放大的上界) & 3.556e-10 \checkmark \\
A & $\kappa$=1e5、n=6 时精确线搜索 CG 的步数与共轭偏差 & 步数 > n 且共轭偏差 \textasciitilde{}1(\allowbreak{}浮点算术破坏共轭性) & 8 步(\allowbreak{}n=\allowbreak{}6)，偏差 0.793，终距 1.25e-10 \checkmark \\
A & CG(Wolfe 不精确线搜索) 迭代数 / n & 远大于 1(\allowbreak{}失去共轭性的代价) & 72 /\allowbreak{} 4 =\allowbreak{} 18.0(\allowbreak{}终距 1.55e-05) \checkmark \\
A & 实测特征值相对误差 max |$\lambda$-$\lambda$|/$\lambda$ & $\le$ 1000u $\approx$ 1.1e-13(\allowbreak{}$\kappa$=\allowbreak{}1e3 时 $\approx$ $\kappa$$\cdot$u) & 1.255e-14(\allowbreak{}$\kappa$=\allowbreak{}10: 实测 10.0000000000；$\kappa$=\allowbreak{}100: 实测 100.0000000000；$\kappa$=\allowbreak{}1000: 实测 1000.0000000000) \checkmark \\
A & 最优解残差 $\|$A x* - b$\|$/$\|$b$\|$ & $\le$ 100u $\approx$ 1.1e-14 & 0.000e+00 \checkmark \\
A & Rosenbrock 的二次窗口 & enter\_iter=\allowbreak{}None、\allowbreak{}C=\allowbreak{}None(\allowbreak{}稳定后缀 < 2 点) & 窗口 [](\allowbreak{}0 点)，C=\allowbreak{}None，牛顿 22 步 \checkmark \\
A & 四次型的二次窗口 & $\ge$2 点稳定后缀且 C 有限 & enter\_iter=\allowbreak{}2，C=\allowbreak{}9.905e-05，窗口 [9.905e-05 6.835e-05] \checkmark \\
A & 逃逸阈值相对误差 max |$\delta$*\_\allowbreak{}meas - $\delta$*\_\allowbreak{}theory|/$\delta$*\_\allowbreak{}theory & $\le$ 1e-6 & 5.301e-10(\allowbreak{}$\lambda$-=\allowbreak{}-0.5: 实测 7.396886e-22 /\allowbreak{} 理论 7.396886e-22(\allowbreak{}相对差 5.30e-10)；$\lambda$-=\allowbreak{}-0.05: 实测 2.866052e-03 /\allowbreak{} 理论 2.866052e-03(\allowbreak{}相对差 3.25e-11)) \checkmark \\
A & 45$^\circ$ 方向阈值 / u\_\allowbreak{}min 方向阈值(对 1/|<u\_\allowbreak{}min,d>|) & =\allowbreak{} 1.414214 & 1.414214(\allowbreak{}重合度 0.707107) \checkmark \\
A & 正交方向的逃逸结果 & escaped=\allowbreak{}False、\allowbreak{}threshold=\allowbreak{}inf、\allowbreak{}重合度=\allowbreak{}0 & escaped=\allowbreak{}False，threshold=\allowbreak{}inf，重合度=\allowbreak{}0.000e+00，最大试探 $\delta$=\allowbreak{}1.100e+06 \checkmark \\
A & 同一问题上沿 u\_\allowbreak{}min 的阈值 & 有限(\allowbreak{}$\approx$1.4e-3) & 2.866e-03(\allowbreak{}escaped=\allowbreak{}True) \checkmark \\
A & 与参考实现终点的最大距离 max $\|$x\_\allowbreak{}ours - x\_\allowbreak{}ref$\|$ & $\le$ 1e-8 & 8.937e-15(\allowbreak{}newton: $\|$$\Delta$x$\|$=\allowbreak{}0.00e+00(\allowbreak{}22 步，f=\allowbreak{}0.00e+00)；bfgs: $\|$$\Delta$x$\|$=\allowbreak{}8.94e-15(\allowbreak{}36 步，f=\allowbreak{}1.60e-29)；trust\_region: $\|$$\Delta$x$\|$=\allowbreak{}0.00e+00(\allowbreak{}22 步，f=\allowbreak{}0.00e+00)) \checkmark \\
A & 参考实现终点与解析极小点 (1,1) 的距离 & $\le$ 1e-12(\allowbreak{}参考值本身可信) & 0.000e+00(\allowbreak{}f=\allowbreak{}0.000e+00) \checkmark \\
A & 实测压缩因子与 max|1-$\alpha$$\lambda$| 的最大相对差 & $\le$ 1e-4 & 1.132e-16(\allowbreak{}16 个步长，$\alpha$/\allowbreak{}$\alpha$* $\in$ [0.7,\allowbreak{}1.0]) \checkmark \\
A & 实测最优步长 $\alpha$\_\allowbreak{}argmin/$\alpha$* & =\allowbreak{} 1(\allowbreak{}网格分辨率 0.02) & 1.0000(\allowbreak{}$\alpha$*=\allowbreak{}1.980198e-02，最小实测因子 0.980198) \checkmark \\
A & $\alpha$* 处的实测因子 vs ($\kappa$-1)/($\kappa$+1) & 0.980198 & 0.980198 \checkmark \\
A & 实测迭代数与理论 ceil(ln$\varepsilon$/ln$\rho$) 的最大相对差 & $\le$ 1e-3 & 0.000e+00(\allowbreak{}$\kappa$=\allowbreak{}10000:92104；$\kappa$=\allowbreak{}3000:27632；$\kappa$=\allowbreak{}1000:9211；$\kappa$=\allowbreak{}300:2764；$\kappa$=\allowbreak{}100:922；$\kappa$=\allowbreak{}30:277；$\kappa$=\allowbreak{}10:92) \checkmark \\
A & N($\kappa$) 的幂律指数 & 1 & 0.99992 \checkmark \\
A & 拟合常数 C = N/$\kappa$ 与 1/2ln(1/$\varepsilon$) 的比 & $\approx$ 1(\allowbreak{}1/\allowbreak{}2ln(\allowbreak{}1/\allowbreak{}$\varepsilon$)=\allowbreak{}9.2103) & C=\allowbreak{}9.21410，比 1.00041 \checkmark \\
A & 网格值相对误差 max & =\allowbreak{} 0(\allowbreak{}同一公式，逐点求值) & 0.000e+00(\allowbreak{}17$\times$17 网格，抽样 25 点) \checkmark \\
A & 网格坐标范围 & X $\in$ [-2,\allowbreak{}2]、\allowbreak{}Y $\in$ [-1,\allowbreak{}3]，shape=\allowbreak{}(\allowbreak{}17,\allowbreak{}17) & X$\in$[-2.0,\allowbreak{}2.0]，Y$\in$[-1.0,\allowbreak{}3.0]，shape=\allowbreak{}(\allowbreak{}17,\allowbreak{} 17) \checkmark \\
A & B 精确满足 / B 不满足的残差 & 0 与 0.316(\allowbreak{}=\allowbreak{}$\|$(\allowbreak{}1,\allowbreak{}2)$\|$/\allowbreak{}$\|$(\allowbreak{}2,\allowbreak{}6)$\|$) & 0.000e+00 与 0.353553 \checkmark \\
A & y=0 时的返回值 & inf(\allowbreak{}无定义，不许返回 0) & inf \checkmark \\
C & GD 收敛因子($\kappa$=100) & 0.980198 & 0.980198 \checkmark \\
C & GD 收敛因子($\kappa$=1000) & 0.998002 & 0.998002 \checkmark \\
C & GD 收敛因子($\kappa$=10000) & 0.9998 & 0.9998 \checkmark \\
C & 实测收敛因子与 ($\kappa$-1)/($\kappa$+1) 的最大相对差 & $\le$ 1e-5 & 3.331e-16(\allowbreak{}$\kappa$=\allowbreak{}100: 实测 0.980198 /\allowbreak{} 理论 0.980198；$\kappa$=\allowbreak{}1000: 实测 0.998002 /\allowbreak{} 理论 0.998002；$\kappa$=\allowbreak{}10000: 实测 0.999800 /\allowbreak{} 理论 0.999800) \checkmark \\
C & 动量实测速率与 ($\surd$$\kappa$-1)/($\surd$$\kappa$+1) 的最大相对差 & $\le$ 0.02 & 0.0086(\allowbreak{}$\kappa$=\allowbreak{}100: $\rho$ 实测 0.825219/\allowbreak{}理论 0.818182，N 比 8.06；$\kappa$=\allowbreak{}1000: $\rho$ 实测 0.940795/\allowbreak{}理论 0.938693，N 比 24.39；$\kappa$=\allowbreak{}10000: $\rho$ 实测 0.980840/\allowbreak{}理论 0.980198，N 比 74.13) \checkmark \\
C & N\_\allowbreak{}gd/N\_\allowbreak{}momentum $\div$ $\surd$$\kappa$ & 实测 0.70--0.85(\allowbreak{}**不是** 1：速率是渐近的) & [0.8056 0.7713 0.7413] \checkmark \\
C & 最大可收敛步长 / $\alpha$* & 1.00942 & 1.00986 \checkmark \\
C & 临界步长与 2/L 的比值 & 1(\allowbreak{}$\alpha$\_thr $\to$ 2/\allowbreak{}L，见 theory 里的 0.1\textasciicircum{}(\allowbreak{}1/\allowbreak{}N) 修正) & $\alpha$\_thr/\allowbreak{}$\alpha$*=\allowbreak{}1.009860674，$\alpha$*/\allowbreak{}(\allowbreak{}2/\allowbreak{}L)=\allowbreak{}0.990099010 \checkmark \\
C & 二分区间宽度 / 谓词调用次数 & 区间宽度 $\le$ 1e-9(\allowbreak{}相对 1) & 9.313e-10(\allowbreak{}30 次二分，32 次谓词调用) \checkmark \\
C & $\kappa$\_\allowbreak{}max(N 步预算) & 217.147 & 217.04 \checkmark \\
C & $\kappa$\_\allowbreak{}max 实测 / 2N$\cdot$ln(1/$\varepsilon$)\textasciicircum{}(-1) & 1(\allowbreak{}相对 5e-3) & 217.0402 /\allowbreak{} 217.1472 =\allowbreak{} 0.999507 \checkmark \\
C & L-BFGS 迭代数随 m 的序列 & 不增(\allowbreak{}m 越大保留的曲率信息越多) & m=\allowbreak{}[1,\allowbreak{} 2,\allowbreak{} 3,\allowbreak{} 5,\allowbreak{} 8,\allowbreak{} 12,\allowbreak{} 20] $\to$ [169,\allowbreak{} 95,\allowbreak{} 34,\allowbreak{} 29,\allowbreak{} 25,\allowbreak{} 22,\allowbreak{} 22](\allowbreak{}完整 BFGS 6 步) \checkmark \\
C & N(m=1)/N(m=20) 与 N(m=20)/N(BFGS) & $\ge$ 5 与 $\le$ 6(\allowbreak{}记忆够用时接近完整 BFGS) & 7.68 与 3.67 \checkmark \\
\end{longtable}


\section{它会怎样骗你（本章版）}

\begin{enumerate}
  \item \textbf{"曲线很陡"不是二次收敛。} 四次型上 $q_k$ 的实测序列是
        $0.1076,\ 0.0214,\ 9.905\times10^{-5},\ 6.835\times10^{-5}$：
        只有最后两步常数（散布 $1.45$），第一步比渐近常数大 $1087$ 倍。
        判据是"$q_k$ 趋于常数"，不是"误差掉得快"。
  \item \textbf{舍入地板之后的 $q$ 爆炸会被读成"超二次收敛"。} 同一算例里 $q$ 从
        $6.8\times10^{-5}$ 一路冲到 $3.35\times10^{122}$；真相是误差已到机器精度、
        转而按每步 $\times u$ 线性下降，$q_k=e_{k+1}/e_k^2$ 必然发散。
  \item \textbf{卡住步长的牛顿法仍被叫作"二次收敛"。} $\alpha_{\text{cap}}=0.5$ 时实测线性因子
        $0.5000$、$\alpha_{\text{cap}}=0.25$ 时 $0.7506$；尾部阶从 $\approx3$ 掉到 $1.000$。
        二次收敛是"步长策略 $+$ 初值位置 $+$ 曲率符号"三件事的属性，不是方法的属性。
  \item \textbf{在 Rosenbrock 上硬报一个二次常数。} 尾部 $q$ 是
        $4.09,\ 1.81,\ 13.9,\ 1.98,\ 18.2$，相邻比值 $9.2$，
        稳定后缀只剩 $1$ 个点：常数必须写 \texttt{null}。
  \item \textbf{把"纯牛顿在 Rosenbrock 上一定发散"写成事实。} 实测（$\varepsilon=10^{-14}$）
        $(-1.2,1)/(5,25)/(-3,10)$ 分别 $7/3/7$ 步收敛，这个经典反例在本章算例上\textbf{没有复现}，
        于是不写；真要演示牛顿法的麻烦，就用"卡步长退化为线性"或者鞍点那两条。
  \item \textbf{"BFGS 满足割线方程"只报更新之后。} 更新后 $3.05\times10^{-15}$、
        更新前 $0.999$（第一次迭代前 $0.9989$），差 $12$ 个数量级。
        割线方程是更新公式的\textbf{目标}，不是恒等式。
  \item \textbf{"L-BFGS 满足割线方程"不说是哪一对。} $m=2$ 时最新一对 $0$、上一对 $0.234$；
        $m=5$ 时最老一对 $1.688$。只精确满足最近 $m$ 对。
  \item \textbf{"动量把迭代数降到 $\sqrt\kappa$ 分之一"。} 速率对得上（相对差 $\le0.009$），
        迭代数只降到 $0.81/0.77/0.74$ 倍的 $\sqrt\kappa$；
        差的那两成是瞬态因子 $(1+k)$。
  \item \textbf{"CG 总是 $n$ 步收敛"。} 精确线搜索下 $(n,\kappa)=(4,10^3)$ 是 $4$ 步，
        换成 Wolfe 线搜索 $72$ 步；$\kappa=10^5$ 时即使精确线搜索也要 $8$ 步、
        共轭偏差 $0.995$。
  \item \textbf{"随机扰动逃出鞍点"当成方法有效。} 沿 $u_{\min}$ 阈值 $2.866\times10^{-3}$、
        沿 $45^\circ$ 是 $4.053\times10^{-3}$（比 $1.414=1/\text{重合度}$）；
        与 $u_{\min}$ 正交时放大 $10^{6}$ 倍也逃不掉，阈值是"没测到"不是 $0$。
  \item \textbf{三种停机判据给出三个结论。} 同一个 $\kappa=1000$ 的二次型，
        梯度判据要 $9211$ 步（与理论完全相等），而
        $\kappa_{\max}\approx2N/\ln(1/\varepsilon)$ 在 $N=2000$ 时给出实测阈值 $217.04$
        （理论 $217.15$）——先说是哪个判据，再谈"收敛了没有"。
  \item \textbf{发散被写成 \texttt{max\_iter}}。$\kappa=10^4$、\texttt{step\_scale}$=1.6$ 时
        跑满 $400$ 步才长到 $10^{137}$，够不到旧的绝对守卫 $10^{150}$。
        发散判据必须是相对的：比初值放大 $1/u\approx4.5\times10^{15}$ 倍就算发散。
  \item \textbf{用二次型演示"信赖域半径在自适应"。} 精确模型下 $\rho\equiv1.000000$、
        零次拒绝、半径只增不减（$1\to2\to2$）；
        换 Rosenbrock 才有 $3$ 次拒绝、$\rho=-4.25$、半径在 $0.125$ 与 $1.0$ 之间来回。
  \item \textbf{阈值参数 $t_2$ 取到区间外。} $t_2\ge1$ 时由于 $\rho\le1$，
        半径永不扩张（实测 $t_2=1.5$ 时区间 $[0.25,1.0]$、$\Delta_0=1.0$）；
        $t_2\le t_1$ 时收缩与扩张的判据互相打架。
  \item \textbf{把占位值当真值画。} 对数轴上的 $0$、模型下降量 $\le0$ 时无定义的 $\rho_k$、
        轨迹被拒绝的步，都必须有掩码、地板和文字说明三件套；
        否则图上"少一段"会被读成"一直是那个值"。
\end{enumerate}
